\section{Results and Discussion}
\label{sec:results}

Table~\ref{tab:main-hazard-results} shows the aligned full-population hazard
results. The full DynamicSTGNN obtains the strongest performance across all
reported ranking and probability metrics. Compared with the no-graph neural
baseline, the full graph model improves AUC by 0.0100 and AP by 0.0112. Compared
with the tabular baseline, it improves AUC by 0.0130 and AP by 0.0181 on the
same aligned test cells. The edge ablation also shows a monotonic pattern:
no graph \(<\) proximity \(<\) proximity+combat \(<\) full graph.

\begin{table}[t]
\centering
\caption{Aligned test-set hazard prediction results. All rows use the same
134,245 target-defined team-phase cells.}
\label{tab:main-hazard-results}
\small
\begin{tabular}{lrrrr}
\toprule
Model & AUC & AP & Acc. & BCE \\
\midrule
Tabular HGB & 0.7651 & 0.6626 & 0.7483 & 0.5140 \\
DynamicSTGNN, no graph & 0.7681 & 0.6695 & 0.7518 & 0.5090 \\
DynamicSTGNN, proximity & 0.7741 & 0.6762 & 0.7545 & 0.5046 \\
DynamicSTGNN, proximity+combat & 0.7755 & 0.6779 & 0.7556 & 0.5036 \\
DynamicSTGNN, full graph & \textbf{0.7782} & \textbf{0.6808} & \textbf{0.7568} & \textbf{0.5016} \\
\bottomrule
\end{tabular}
\end{table}

The paired tests in Table~\ref{tab:paired-tests} show that these differences
are not driven by a small number of matches. The full graph model improves over
the no-graph model on match-level BCE, AUC, and AP with one-sided Wilcoxon
\(p < 10^{-22}\) in all three cases. The improvements over the tabular baseline
are also significant under the same match-paired test. The full graph model
also improves over proximity-only and proximity+combat variants, which indicates
that co-pressure adds signal beyond the other two edge channels. The advantage
over the no-graph control is positive at every cut phase (+0.007 to +0.025
AUC) and is largest in the sparse late phases, and it also holds within every
mode---six of six, from +0.008 (third-person squad) to +0.020 AUC---so the
gain is not a mode shortcut. This supports
the claim that the observed dynamic team relations contain useful signal beyond
the team-phase feature sequence alone.

\begin{table}[t]
\centering
\caption{Match-paired Wilcoxon tests on 2,182 test matches. Positive
improvement means the full graph model is better.}
\label{tab:paired-tests}
\small
\begin{tabular}{llrr}
\toprule
Comparison & Metric & Mean impr. & \(p\) \\
\midrule
Full vs. no graph & BCE & 0.0071 & \(1.63{\times}10^{-47}\) \\
Full vs. no graph & AUC & 0.0080 & \(1.17{\times}10^{-30}\) \\
Full vs. no graph & AP  & 0.0090 & \(3.12{\times}10^{-23}\) \\
Full vs. proximity & BCE & 0.0030 & \(2.07{\times}10^{-15}\) \\
Full vs. proximity & AUC & 0.0039 & \(4.61{\times}10^{-10}\) \\
Full vs. proximity & AP  & 0.0051 & \(1.89{\times}10^{-10}\) \\
Full vs. prox.+combat & BCE & 0.0019 & \(1.89{\times}10^{-9}\) \\
Full vs. prox.+combat & AUC & 0.0022 & \(1.69{\times}10^{-5}\) \\
Full vs. prox.+combat & AP  & 0.0028 & \(7.71{\times}10^{-6}\) \\
Full vs. tabular & BCE & 0.0120 & \(3.30{\times}10^{-74}\) \\
Full vs. tabular & AUC & 0.0113 & \(2.60{\times}10^{-37}\) \\
Full vs. tabular & AP  & 0.0154 & \(9.28{\times}10^{-36}\) \\
\bottomrule
\end{tabular}
\end{table}

Calibration follows the same pattern. As shown in
Table~\ref{tab:calibration-results}, the full graph model has the lowest Brier
score, BCE, and ECE. Its mean predicted hazard probability, 0.3260, is close to
the empirical positive rate of 0.3245. The tabular baseline is still a strong
model, but it is the least calibrated in this aligned evaluation; its ECE is
roughly five times that of the full graph model.

\begin{table}[t]
\centering
\caption{Calibration diagnostics on aligned test cells. ECE uses 10
equal-width probability bins.}
\label{tab:calibration-results}
\small
\begin{tabular}{lrrrr}
\toprule
Model & Brier & BCE & ECE & Mean prob. \\
\midrule
Tabular HGB & 0.1708 & 0.5140 & 0.0075 & 0.3223 \\
DynamicSTGNN, no graph & 0.1692 & 0.5090 & 0.0039 & 0.3264 \\
DynamicSTGNN, proximity & 0.1674 & 0.5046 & 0.0023 & 0.3253 \\
DynamicSTGNN, proximity+combat & 0.1670 & 0.5036 & 0.0028 & 0.3266 \\
DynamicSTGNN, full graph & \textbf{0.1662} & \textbf{0.5016} & \textbf{0.0015} & 0.3260 \\
\bottomrule
\end{tabular}
\end{table}

Table~\ref{tab:feature-ablation} answers part (i) of RQ2 at the
representation level. Restricted to the five-variable score summary, the same
learner reaches pooled AUC 0.649; adding the raw quantities behind those
scores, conventional aggregates, and the STDB movement and exposure features
raises pooled AUC to 0.765. For significance testing, we compute AUC and AP
within each test match and then average the paired improvements. Under this
match-paired calculation, the full feature set gains +0.112 AUC and +0.170 AP
over the score-summary representation (one-sided \(p\) below numerical
precision; rank-biserial 0.96 for AUC and 0.97 for AP). The STDB-specific
features contribute +0.012 AUC and +0.014 AP over conventional aggregates
(\(p<10^{-62}\)), and the sixteen STDB spatio-temporal features alone beat
the five-variable score summary by +0.057 AUC. Appending the derived scores
to the full set adds nothing (\(-0.0003\) AUC, \(p=0.75\)).

\begin{table}[t]
\centering
\caption{Representation ablation: the same gradient-boosting pipeline on the
same aligned test cells, varying only the feature set (\(d\) features).}
\label{tab:feature-ablation}
\small
\begin{tabular}{lrrrr}
\toprule
Feature set & \(d\) & AUC & AP & BCE \\
\midrule
Prior-style score summary & 5 & 0.6490 & 0.4547 & 0.5993 \\
Score raw inputs & 12 & 0.7184 & 0.5703 & 0.5584 \\
STDB spatio-temporal only & 16 & 0.7124 & 0.5895 & 0.5529 \\
Conventional aggregates & 57 & 0.7522 & 0.6476 & 0.5242 \\
Conventional + STDB spatio-temporal & 72 & \textbf{0.7652} & \textbf{0.6626} & \textbf{0.5141} \\
Above + prior scores & 76 & 0.7649 & 0.6624 & 0.5141 \\
\bottomrule
\end{tabular}
\vspace{2pt}
\begin{minipage}{0.96\linewidth}
\footnotesize
\textit{Note.} Score raw inputs are 12 variables covering phase index, combat
counts, damage, knockdowns, distinct opponents damaged, path length, average
speed, 100m enemy density, item pickups, unique items, and care-package
pickups. The score summary uses phase index plus fight score, active score,
strategy type, and type survival rate derived from those raw inputs.
\end{minipage}
\end{table}

These results close the loop on the two research questions. RQ1 is answered by
construction: 21,824 of 27,491 raw matches survive conversion into
phase-aligned, graph-ready form (Table~\ref{tab:data-funnel}), and every number
above is computed from that representation. Part (i) of RQ2 also receives a
clear answer: feature-only models reach AUC 0.765--0.768 with near-nominal
calibration, so phase-level elimination hazard is substantially learnable from
the team-phase features alone. Table~\ref{tab:feature-ablation} traces this to the representation rather
than the model class: pooled AUC falls from 0.765 to 0.649 when the same
learner is restricted to prior-style strategy scores. Part (ii) is the central finding: dynamic relations between teams add
predictive information beyond those features. The graph edges improve both
ranking metrics and probability quality over the no-graph control. Spatially
and interactionally defined team relations are thus part of the
elimination-risk environment, consistent with the habitat-contraction reading
of Section~\ref{sec:related-work}.

At the same time, the result should be read as predictive rather than causal.
The model shows that observed proximity, combat, and co-pressure relations help
predict next-phase elimination hazard. It does not prove that changing a team's
strategy to alter one of these relations would causally change survival
probability. Shuffling team identities within the observed graphs at
inference time collapses the full model to AUC 0.701 over five repeats
(\(-0.077\)), well below even the no-graph control, so the relational gain
reflects who relates to whom rather than graph regularization. Within the
edge ablation, proximity contributes the largest single step
(+0.0060 AUC over no graph), and the full model remains significantly better
than both partial variants under the match-paired tests; repeated-seed
confirmation on the full population is left as future work.
